\documentclass[12pt, answers]{exam}
\usepackage{tikz}
\usepackage{amsmath}
\usepackage{amsfonts}
\usepackage{algpseudocodex}
\usepackage{algorithm}
\usepackage{amssymb}
\usepackage{graphicx}
\usepackage{multirow}
\usepackage{enumitem}
\usepackage{subfiles}
\usepackage[most]{tcolorbox} 
\usepackage{changepage} % for adjustwidth
\usepackage{xepersian}

\settextfont[
    Path = ../fonts/,
    Scale = 1.0,
    BoldFont = *Bd,
    ItalicFont = *It,
    BoldItalicFont = *BdIt,
    Extension = .ttf
]{XB Yas}

\renewenvironment{solution}{%
  \par
  \begin{tcolorbox}[
    breakable,                
    colback=white,
    colframe=black,
    arc=1pt,
    boxrule=0.6pt,
    left=4mm, right=4mm,
    top=2mm, bottom=2mm,
    before skip=10pt,
    after skip=10pt,
    segmentation style={solid, black}, 
    before upper={\noindent\textbf{\color{blue}پاسخ:}\par\vspace{0.5em}}
  ]
 }
{%
  \end{tcolorbox}
  \vspace{1em}
}


\algrenewcommand{\algorithmicrequire}{\textbf{Input:}}
\algrenewcommand{\algorithmicensure}{\textbf{Output:}}

\definecolor{codegreen}{rgb}{0,0.6,0}
\definecolor{codegray}{rgb}{0.5,0.5,0.5}
\definecolor{codepurple}{rgb}{0.58,0,0.82}
\definecolor{backcolour}{rgb}{0.95,0.95,0.92}

\lstdefinestyle{mystyle}{
    backgroundcolor=\color{backcolour},   
    commentstyle=\color{codegreen},
    keywordstyle=\color{magenta},
    numberstyle=\tiny\color{codegray},
    stringstyle=\color{codepurple},
    basicstyle=\ttfamily\footnotesize,
    breakatwhitespace=false,         
    breaklines=true,                 
    captionpos=b,                    
    keepspaces=true,                 
    numbers=left,                    
    numbersep=5pt,                  
    showspaces=false,                
    showstringspaces=false,
    showtabs=false,                  
    tabsize=2
}

\lstset{style=mystyle}

% Wide solution environment (extends into the question label area)
% Usage: replace \begin{solution} ... \end{solution} with \begin{widesolution} ... \end{widesolution}
% Keeps the same optional title argument as exam's solution: \begin{widesolution}[عنوان]
\newenvironment{widesolution}[1][]{%
    \begin{solution}[#1]%
    % Outdent left by the current list left margin so the solution spans full text width
    \begin{adjustwidth}{\dimexpr-\leftmargin\relax}{0pt}%
}{%
    \end{adjustwidth}%
    \end{solution}%
}

\renewcommand{\solutiontitle}{\noindent\textbf{\color{blue}پاسخ:}\par}
\newcommand{\class}{اصول سیستم‌های عامل}
\newcommand{\term}{نیم‌سال اول ۰۴-۰۵}
\newcommand{\college}{دانشکده ریاضی و آمار}
\newcommand{\prof}{استاد: دکتر رفیعی}
\newcommand{\group}{\textcolor{red}{نام دانشجو}}

\begin{document}
%%%%%%%%%%%%%% HDR START 
\pagestyle{head}
\firstpageheader{}{}{}
\runningheader{صفحه \thepage\ از \numpages}{}{\class}
\runningheadrule
\begin{tabular}{p{.7\textwidth} l}
    \multicolumn{2}{c}{\textbf{به نام خدا}}                                            \\
                    & \multirow{2}{*}{\includegraphics[scale=0.2]{../UILOGO.png}} \\
    \textbf{\class} &                                                                  \\
    \textbf{\term}  &                                                                  \\
    \textbf{\prof}  &                                                                  \\
    \textbf{\group} &                                                                  \\
                    & \textbf{\college}
    \\
    % & \textbf{\todaydate} \\
    % & \textbf{\deadline} \\ \\
\end{tabular}\\
\rule[1ex]{\textwidth}{.1pt}

\begin{center}
    \textbf{\lr{Problem Set 1}}
\end{center}
%%%%%%%%%%%%%% HDR END 


\begin{questions}
    \pointpoints{نمره}{نمره}

    %%%%%%%%%%%%%% 1 START
    \question  صحیح یا غلط بودن گزاره‌های زیر را تعیین کنید و در صورت غلط بودن، توضیح دهید چرا غلط است.
\begin{parts}

    \part قالب داده و محل قرارگیری آن در ناحیه حافظه اشتراکی به طور مستقیم توسط سیستم عامل کنترل و مدیریت می‌شود.

    \part به طور معمول، ناحیه حافظه اشتراکی در فضای آدرس فرآیندی که ارتباط را آغاز می‌کند (فرآیند ارتباطی) ایجاد می‌شود.

    \part در مدل «بافر نامحدود»
    \lr{(Unbounded Buffer)}،
    فرآیند مصرف‌کننده هرگز در حالت انتظار قرار نمی‌گیرد.

    \part در آدرس‌دهی مستقیم «نامتقارن»
    \lr{(Asymmetric)}
    ، گیرنده برای دریافت پیام باید نام فرستنده را بداند.

    \part در بافر با «ظرفیت صفر» (سیستم پیام بدون بافر)، فرستنده باید تا زمانی که گیرنده پیام را دریافت کند، بلاک شود.

    \part وب‌سرورها برای مدیریت درخواست‌های هم‌زمان کاربران، معمولاً برای هر درخواست یک «فرآیند» جدید ایجاد می‌کنند زیرا سربار کمتری دارد.

    \part در مدل «چند به یک»
    \lr{(Many-to-One)}،
    اگر یکی از نخ‌های سطح کاربر یک فراخوان سیستمی مسدودکننده
    \\ \lr{(Blocking System Call)
    } انجام دهد، کل فرآیند بلاک می‌شود.

    \part مدیریت نخ‌های سطح کاربر
    \lr{(User Threads)
    } نیازمند پشتیبانی مستقیم سیستم‌عامل (کرنل) است.

    \part در مدل «چند به یک»، نخ‌های سطح کاربر می‌توانند به صورت موازی روی سیستم‌های چند هسته‌ای اجرا شوند.

    \part در استراتژی «نخ ضمنی»
    \lr{(Implicit Threading)
    }، وظیفه ایجاد و مدیریت نخ‌ها از برنامه‌نویس به سیستم‌عامل (کرنل) منتقل می‌شود.

    \part
    در مدل‌های چند به چند، اگر نخ سطح کرنل متصل به یک LWP بلاک شود (
    مثلاً برای
    I/O
    )، نخ سطح کاربر متصل به آن LWP همچنان می‌تواند به اجرا ادامه دهد.

\end{parts}
    \begin{solution}
        \begin{parts}
            \part غلط. قالب داده و محل آن توسط فرآیندهای همکار کنترل می‌شود، نه سیستم عامل.
            \part صحیح.
            \part غلط. حتی در بافر نامحدود، اگر بافر خالی باشد (انتظار برای اضافه شدن عنصر جدید)، مصرف‌کننده ممکن است منتظر بماند.
            \part غلط. در طرح نامتقارن، گیرنده نیازی به نام فرستنده ندارد؛ بلکه می‌تواند از هر شناسه‌ای (id) دریافت کند.
            \part صحیح.
            \part غلط. معمولاً از نخ‌ها استفاده می‌شود زیرا سربار کمتری دارند.
            \part صحیح.
            \part غلط. مدیریت آن‌ها بدون پشتیبانی سیستم‌عامل و در فضای کاربر انجام می‌شود.
            \part غلط. چون تنها یک نخ کرنل وجود دارد، تنها یک نخ در یک زمان قابل اجراست و امکان اجرای موازی وجود ندارد.
            \part غلط. این وظیفه به کامپایلرها و کتابخانه‌ها منتقل می‌شود.
            \part غلط. اگر نخ کرنل بلاک شود، LWP و متعاقباً نخ کاربر متصل به آن نیز بلاک خواهند شد.
        \end{parts}


    \end{solution}
    %%%%%%%%%%%%%% 1 END

    \pagebreak

    %%%%%%%%%%%%%% 2 START
    \question[12] به سوالات زیر پاسخ مناسب دهید.

\begin{parts}

    \part ساختار داده مناسب برای پیاده‌سازی بافر محدود
    \lr{(Bounded Buffer)}
    چیست؟

    \part دو فرآیند برای برقراری ارتباط از طریق حافظه اشتراکی، باید چه محدودیت خاصی که توسط سیستم‌عامل اعمال می‌شود را حذف کنند؟

    \part فرآیندها چه تضمینی را باید در مورد نوشتن در مکان حافظه اشتراکی ارائه دهند؟

    \part    پیوند ارتباطی به چه روش‌هایی می‌تواند به صورت فیزیکی پیاده‌سازی شود؟

    \part هنگامی که مالکیت صندوق پستی با سیستم عامل است، سیستم عامل چه سه مکانیزمی را باید فراهم کند؟

    \part مزیت «پاسخگویی سریع» در برنامه‌های چندنخی به چه معناست؟

    \part یک نخ (Thread) به عنوان واحد پایه استفاده از
    \lr{CPU}،
    شامل چه ۴ جزء اختصاصی است؟

    \part چه منابعی بین تمام نخ‌های یک فرآیند مشترک هستند؟ (۳ مورد نام ببرید)

    \part چهار مزیت اصلی برنامه‌نویسی چندنخی که در اسلایدها ذکر شده است را نام ببرید.

    \part پنج چالش اصلی برنامه‌نویسی در سیستم‌های چند‌هسته‌ای را نام ببرید.

    \part چرا مدل «یک به یک» هم‌روندی (Concurrency) بیشتری نسبت به مدل «چند به یک» دارد؟

    \part دو مشکل اصلی ایجاد نخ جدید برای هر درخواست (بدون استفاده از استخر نخ) را نام ببرید.

    \part تفاوت اصلی TLS با داده‌های ایستا
    \lr{(Static Data)}
    چیست؟
\end{parts}

    \begin{solution}
\begin{parts}
    
    \part صف حلقوی \lr{(Circular Queue)}.
    \part  آن‌ها باید توافق کنند که محدودیت سیستم عامل را که در حالت عادی مانع دسترسی یک فرآیند به فضای آدرس فرآیند دیگر می‌شود، حذف کنند.
    \part{۱۴.} فرآیندها باید ضمانت کنند که به طور همزمان روی مکان یکسانی نمی‌نویسند.
    \part{۱۵.} از طریق حافظه اشتراکی، گذرگاه سخت‌افزاری، یا شبکه.
    \part{۱۶.} ۱. امکان ایجاد صندوق پستی جدید. ۲. امکان ارسال و دریافت پیام از طریق صندوق پستی. ۳. امکان حذف صندوق پستی.
    \part{۱۷.} به این معنی است که اگر بخشی از برنامه بلاک شود، سایر بخش‌ها می‌توانند به اجرا ادامه دهند.
    \part{۱۸.} ۱. شناسه نخ \lr{(Thread ID)}، ۲. شمارنده برنامه \lr{(Program Counter)}، ۳. مجموعه ثبات \lr{(Register Set)}، ۴. پشته \lr{(Stack)}.
    \part{۱۹.} ۱. بخش کد، ۲. بخش داده، ۳. منابع سیستم‌عامل (مانند فایل‌های باز).
    \part{۲۰.} ۱. پاسخگویی سریع، ۲. اشتراک‌گذاری منابع، ۳. صرفه اقتصادی، ۴. مقیاس‌پذیری.
    \part{۲۱.} ۱. شناسایی وظایف، ۲. تعادل، ۳. جداسازی داده، ۴. وابستگی داده، ۵. تست و اشکال‌زدایی.
    \part{۲۲.} زیرا بلاک شدن یک نخ در سطح کاربر باعث بلاک شدن کل فرآیند نمی‌شود و سایر نخ‌ها می‌توانند به کار خود ادامه دهند.
    \part{۲۳.} ۱. سربار زمانی ایجاد نخ و حذف آن پس از اتمام کار. ۲. عدم وجود محدودیت برای تعداد نخ‌های فعال که می‌تواند باعث هدررفت منابع شود.
    \part{۲۴.} داده‌های ایستا معمولاً بین همه نخ‌ها مشترک هستند، اما داده‌های TLS برای هر نخ یکتا و اختصاصی هستند.
\end{parts}
    \end{solution}
    %%%%%%%%%%%%%% 2 END

    \pagebreak

    %%%%%%%%%%%%%% 3 START 
    \question[14] 
در یک سیستم‌عامل ، دو تسک از \lr{Mailbox} برای ارسال پیام استفاده می‌کنند.  
وقتی تعداد پیام‌ها زیاد می‌شود، سیستم دچار تأخیر و از دست رفتن Deadline می‌شود.  
توضیح دهید چرا این اتفاق می‌افتد و چه روش‌هایی برای کاهش این مشکل وجود دارد.



    \begin{solution}
        \textbf{دلایل رخ‌دادن:}
        \begin{itemize}
            \item Mailbox ظرفیت محدود دارد و وقتی پر شود، تسک فرستنده بلوکه می‌شود و زمان پاسخ بالا می‌رود.
            \item ارسال و دریافت پیام شامل کپی‌کردن داده و \lr{Context Switch} است که در بار سنگین باعث افزایش تأخیر می‌شود.
            \item اگر Consumer اولویت پایین داشته باشد، خالی‌شدن Mailbox دیر انجام می‌شود و عملاً \lr{Priority Inversion} رخ می‌دهد.
        \end{itemize}

        \textbf{راه‌حل‌ها:}
        \begin{itemize}
            \item بالا بردن اولویت \lr{Consumer}.
            \item ارسال اشاره‌گر به داده به‌جای کپی کامل پیام.
            \item استفاده از حالت غیرمسدودکننده یا \lr{Timeout}.
            \item افزایش ظرفیت Mailbox براساس نرخ واقعی تولید و مصرف.
        \end{itemize}
    \end{solution}
    %%%%%%%%%%%%%% 3 END 

    \pagebreak

    %%%%%%%%%%%%%% 4 START
    \question[12] 
در یک برنامه چندنخی، یک \lr{Thread Pool} با ۸ نخ وجود دارد و سه نوع کار اجرا می‌شود:  
کارهای سنگین 
CPU
، کارهای 
I/O
، و کارهای سبک با اهمیت زمانی.  
گاهی کارهای سبک دیر اجرا می‌شوند و CPU هم بعضی وقت‌ها بی‌استفاده می‌ماند.  
توضیح دهید چرا چنین وضعی پیش می‌آید و چه روش‌هایی می‌تواند آن را بهتر کند.


    \begin{solution}
        \textbf{دلایل رخ‌دادن:}
        \begin{itemize}
            \item صف مشترک باعث \lr{Head-of-Line Blocking} می‌شود و کارهای سبک پشت کارهای سنگین گیر می‌کنند.
            \item کارهای CPU-bound نخ‌ها را برای مدت زیاد اشغال می‌کنند و باعث اشباع Pool می‌شوند.
            \item کارهای I/O-bound نخ را هنگام انتظار بلوکه می‌کنند و CPU گاهی بیکار می‌ماند.
        \end{itemize}

        \textbf{راه‌حل‌ها:}
        \begin{itemize}
            \item جداکردن صف‌ها و ساختن Thread Poolهای مجزا برای انواع کارها.
            \item استفاده از \lr{Work-Stealing} برای کاهش انسداد صف.
            \item استفاده از \lr{Priority Queue} برای اجرای سریع کارهای سبک و حساس.
        \end{itemize}

    \end{solution}
    %%%%%%%%%%%%%% 4 END

    \pagebreak

    %%%%%%%%%%%%%% 5 START
    
\question[10]  دو برنامه همکار میخواهند پیامی را با استفاده از روش اشتراک حافظه بین یکدیگر انتقال دهند.
مشکلات این دو برنامه را یافته و برای هریک راه‌حلی ارائه دهید.

\begin{latin}
    \begin{lstlisting}[language=C]
// Progam 1
#include <stdio.h>
#include <unistd.h>
#include <sys/shm.h>
#include <string.h>

int main()
{
    int shmid;
    void *shared_memory;
    char buffer[1024];
    shmid = shmget((key_t)2345, 512, 0666);
    shared_memory = shmat(shmid, NULL, 0);
    printf("Enter some data to write to shared memory\n");
    read(0, buffer, 1024);
    strcpy(shared_memory, buffer);
    shmdt(shared_memory);
    return 0;
}       
\end{lstlisting}
\end{latin}
\begin{latin}
    \begin{lstlisting}[language=C]
// Progam 2
#include <stdio.h>
#include <sys/shm.h>
#include <string.h>

int main()
{
    int shmid;
    void *shared_memory;
    shmid = shmget((key_t)1234, 512, 0666);
    shared_memory = shmat(shmid, NULL, 0);
    printf("Data read from shared memory: %s\n", shared_memory);
    shmdt(shared_memory);
    return 0;
}    
\end{lstlisting}
\end{latin}


    \begin{solution}

    \end{solution}
    %%%%%%%%%%%%%% 5 END

    \pagebreak

    %%%%%%%%%%%%%%% 6 START
    \question[20] 
سه برنامه مستقل می‌خواهند با استفاده از مکانیزم ارسال پیام
\lr{(Message Passing)}
با یکدیگر ارتباط برقرار کنند:

\begin{itemize}
    \item \textbf{برنامه A:}
          اعداد صحیح را از ورودی دریافت کرده و به صف پیام ارسال می‌کند.
          با ورود عدد \texttt{0} برنامه متوقف می‌شود.

    \item \textbf{برنامه B:}
          فقط اعداد \textbf{زوج} را از صف دریافت کرده و مجموع آن‌ها را چاپ می‌کند.

    \item \textbf{برنامه C:}
          فقط اعداد \textbf{فرد} را از صف دریافت کرده و تعداد آن‌ها را چاپ می‌کند.
\end{itemize}

% کد ناقص برنامه A در زیر آمده است.
توضیح دهید جه تغییراتی در برنامه A لازم است و برنامه‌های B و C چگونه باید پیام‌ها را دریافت کنند؟


این کد را تکمیل کنید تا برنامه A به درستی کار کند.
\begin{latin}
    \begin{lstlisting}[language=C]
// Program A - Sender
#include <stdio.h>
#include <stdlib.h>
#include <sys/msg.h>

struct Message {
    long int msg_type;
    int number;
};

int main() {
    int msgid = msgget((key_t)5678, 0666 | IPC_CREAT);
    struct Message msg;
    int num;
    
    while (1) {
        printf("Enter a number (0 to exit): ");
        scanf("%d", &num);
        if (num == 0) break;
        
        msg.number = num;
        // ...
        
        msgsnd(msgid, &msg, sizeof(int), 0);
    }
    return 0;
}
\end{lstlisting}
\end{latin}


    \begin{solution}

    \end{solution}
    %%%%%%%%%%%%%%% 6 END

    \pagebreak

    %%%%%%%%%%%%%%% 7 START
    \question[20]
برنامه زیر را در نظر بگیرید و به سوالات پاسخ دهید.

\begin{latin}
    \begin{lstlisting}[language=C]
#include <stdio.h>
#include <pthread.h>
#include <unistd.h>

int counter = 0;

void *thread_func(void *arg) {
    int id = *(int *)arg;
    for (int i = 0; i < 3; i++) {
        counter++;
        printf("Thread %d: counter = %d\n", id, counter);
        sleep(1);
    }
    pthread_exit(NULL);
}

int main() {
    pthread_t t1, t2;
    int id1 = 1, id2 = 2;
    
    pthread_create(&t1, NULL, thread_func, &id1);
    pthread_create(&t2, NULL, thread_func, &id2);
    
    pthread_join(t1, NULL);
    pthread_join(t2, NULL);
    
    printf("Final counter = %d\n", counter);
    return 0;
}
\end{lstlisting}
\end{latin}

\begin{parts}
    \part
    اگر خط
    \lr{\texttt{pthread\_join(t1, NULL);}}
    را حذف کنیم، چه اتفاقی می‌افتد؟

    \part
    تفاوت اصلی بین ایجاد یک نخ جدید با \texttt{pthread\_create} و ایجاد یک فرآیند جدید با \texttt{fork} چیست؟
\end{parts}


    \begin{solution}

    \end{solution}
    %%%%%%%%%%%%%%% 7 END


\end{questions}
\end{document}